\emph{Pour chacune des affirmations suivantes, indiquer si elle est vraie ou fausse.\\
Chaque réponse doit être justifiée.\\
Une réponse non justifiée ne rapporte aucun point.}

\bigskip

On considère la suite $\left(u_{n}\right)$ définie par 
\[\left\{\begin{array}{l c l}
u_{0} &= &0\quad \text{et}\\
u_{n + 1} &=& 3 u_{n} + 1\:\: \text{pour tout entier naturel } \:n.
\end{array}\right.\]

\begin{enumerate}
	\item On considère la fonction \texttt{calcul} écrite dans le langage Python qui renvoie la valeur de $u_{n}$.
\begin{center}
\renewcommand\arraystretch{1}
\texttt{
\begin{tabular}{l}
def calcul(n):\\
\quad u = 0\\
\quad for i in range(n):\\
\quad\quad u = 3 * u + 1\\
\quad return u\\
\end{tabular}}
\end{center}

On considère par ailleurs la fonction \texttt{liste} écrite dans le langage Python :

\begin{center}
\renewcommand\arraystretch{1}
\texttt{
		\begin{tabular}{l}
		def liste(n):\\
		\quad l = [ ]\\
		\quad for i in range(n):\\
		\quad\quad l.append( calcul(i) )\\
		\quad return l\\
		\end{tabular}}
\end{center}
		\renewcommand\arraystretch{1.4}

	\textbf{Affirmation 1 :} \og{}l'appel \texttt{liste(6)} renvoie la liste \texttt{[0, 1, 4, 13, 42, 121]}.\fg{}

	\item \textbf{Affirmation 2 :} \og{} pour tout entier naturel $n, \quad u_{n} = \dfrac{1}{2} \times 3^{n}-\dfrac{1}{2}$.\fg{}

	\item \textbf{Affirmation 3 :} \og{} pour tout entier naturel $n$, \quad $u_{n + 1}-u_{n}$ est une puissance de 3.\fg{}
\end{enumerate}
